\section{General results: localisation, hierarchy and the finite-menu geometry}\label{sec:general}

\subsection{Two lemmas on quantifier order}

The first lemma is the abstract form of the separation; the second identifies what aggregation can and cannot lose at a fixed trajectory. The plan-selection gap itself can already occur for static alternative margin vectors.

\begin{lemma}[quantifier inclusion]\label{lem:quant}
Let $\{E_\lambda(z)\}_{\lambda\in\Lambda}$ be any family of plan sets. Then
\[
\Big\{z:\ \textstyle\bigcap_\lambda E_\lambda(z)\neq\varnothing\Big\}\ \subseteq\ \bigcap_\lambda\big\{z:\ E_\lambda(z)\neq\varnothing\big\},
\]
with equality exactly when, for every $z$ in the right-hand set, the family $\{E_\lambda(z)\}_\lambda$ admits a common selector. No topology on the plan space is needed; for a finite menu the check is a finite enumeration.
\end{lemma}

\begin{lemma}[the full cone is lossless at a fixed trajectory]\label{lem:cone}
For $v\in\mathbb R^n$: $v\ge0$ componentwise if and only if $w\cdot v\ge0$ for every $w\in\Wp$.
\end{lemma}

Proofs are in \ref{app:general}. Lemma~\ref{lem:cone} says that an assessor who checks the aggregate floor $w\cdot s\ge0$ at \emph{every} nonnegative weighting, along a \emph{fixed} trajectory, detects exactly the same violations as an assessor who checks each floor separately. Whatever separates the two doctrines must therefore enter through the quantifier over plans in \eqref{eq:central}, not through any blindness of scalarization at a profile.

\subsection{Hierarchy, localisation and weight-family monotonicity}

\begin{proposition}[hierarchy and localisation]\label{prop:hierarchy}
\textup{(i)} For every $w\in\Wp$ and every $z$,
\[
\Etyp(z)\ \subseteq\ E_w(z)\ \subseteq\ \Ephys(z)\ \subseteq\ \Eend(z).
\]
\textup{(ii)} For every $z$, $\Etyp(z)=\bigcap_{w\in\Wp}E_w(z)$; hence $\Vtyp=\{z:\ \bigcap_wE_w(z)\neq\varnothing\}$, and the central inclusion \eqref{eq:central} reads $\Vtyp\subseteq\Vweak$.
\end{proposition}

Part (i) is constraint-set monotonicity under nested safe sets, the standard monotonicity of viability kernels \citep{Aubin1991,Frankowska1989}. Part (ii) is the localisation that gives the paper its subject: the noncompensatory doctrine \emph{is} common-plan robustness over the full cone (Protocols 3 and 4 of Table~\ref{tab:protocols} coincide), so the only difference between the doctrines is the order of the two quantifiers.

\begin{proposition}[weight-family monotonicity]\label{prop:monotone}
For a nonempty weight family $\varnothing\ne W\subseteq\Wp$ let $\mathcal V_W=\bigcap_{w\in W}\mathcal V_w$. Then $\Vtyp\subseteq\mathcal V_W\subseteq\Vphys$, and $W_1\subseteq W_2$ with both families nonempty implies $\mathcal V_{W_2}\subseteq\mathcal V_{W_1}$.
\end{proposition}

Restricting a nonempty weight family -- eliciting a narrower set of admissible weightings -- can only \emph{enlarge} the compensatory accepted set. The nonemptiness condition excludes the vacuous universal test over $W=\varnothing$, whose intersection would be the whole state space. The full cone is thus the strictest compensatory reading, and a separation exhibited at the full cone persists for every nonempty restricted family; the elicitation of weights, however careful, cannot remove it.

\subsection{The finite-menu geometric characterisation}

The mechanism behind the gap is geometric. Call a plan a \emph{margin-class} plan if its coordinate-wise admissibility is governed by a \emph{required-margin vector} $d_a\in\Rp$: plan $a$ is admissible at floor margins $s\ge0$ exactly when $d_a\le s$ componentwise, and it satisfies the destination constraint. A plan whose worst-case tube dips each floor $i$ by at most $d_{a,i}$ and returns is of this class; plans whose admissibility depends on another resource (reserve-financed plans) or that fail the destination are handled separately, as on the witness of Section~\ref{sec:witness}.

A margin vector summarises a path, and the summary must be taken in a declared order. Write $r_{a,i}(t)\ge0$ for the deficit of plan $a$ in coordinate $i$ at time $t\in[0,1]$ along its worst-case tube (the amount by which the tube would breach a floor placed at the initial level), so that $d_{a,i}=\max_t r_{a,i}(t)$.

\begin{lemma}[path-to-margin synchronisation]\label{lem:sync}
For continuous deficit paths and any fixed $w\ge0$, the scalarized path requirement $m_a(w)=\max_{t\in[0,1]}\sum_iw_i\,r_{a,i}(t)$ satisfies $m_a(w)\le w\cdot d_a$, with equality if and only if a single time attains $d_{a,i}$ simultaneously for every coordinate with $w_i>0$.
\end{lemma}
\begin{proof}
Each summand is at most its own maximum, so $\sum_iw_ir_{a,i}(t)\le\sum_iw_id_{a,i}$ for every $t$, which gives the inequality. At a maximiser $t^\ast$ of the continuous scalarized path, equality of the nonnegative weighted sum of these termwise inequalities forces equality in every positive-weight coordinate, i.e.\ $r_{a,i}(t^\ast)=d_{a,i}$ whenever $w_i>0$; conversely a common maximising time gives equality directly.
\end{proof}

The lemma fixes the scope of the geometry that follows. Theorem~\ref{thm:menu} describes the protocol whose scalarized requirement is $w\cdot d_a$ -- within-path maxima taken first, then weighted -- \emph{exactly}; for a protocol that scalarizes first and takes the within-path maximum afterwards it is conservative, the two coinciding precisely when the coordinate dips are synchronised. Neither reading may be substituted for the other without the synchronisation hypothesis. On the witness datum of Section~\ref{sec:witness} each principal plan dips a single coordinate, so the hypothesis holds trivially and the two protocols agree.

\begin{theorem}[finite-menu geometric characterisation]\label{thm:menu}
Let $\mathcal A$ be a finite menu of margin-class plans with required-margin vectors $\mathcal D=\{d_a:a\in\mathcal A\}\subset\Rp$. Assume explicitly that the scalarized requirement for each plan and weight is margin-first, $m_a(w)=w\cdot d_a$; in an exact-tube model this holds, for example, when the coordinate deficits with positive weight attain their maxima at a common time. Then
\begin{enumerate}[label=\textup{(\roman*)},leftmargin=2em]
\item coordinate-wise acceptance is $\{s:\ \exists a,\ d_a\le s\}=\bigcup_{a\in\mathcal A}(d_a+\Rp)$;
\item scalarized (weight-adaptive) acceptance is $\{s:\ \forall w\in\Wp\ \exists a,\ w\cdot d_a\le w\cdot s\}=\conv(\mathcal D)+\Rp$.
\end{enumerate}
Hence the acceptance gap is $\big[\conv(\mathcal D)+\Rp\big]\setminus\bigcup_{a}(d_a+\Rp)$. Replacing the margin menu by $\conv(\mathcal D)$ makes the two abstract acceptance sets coincide. This is a convexification in required-margin space; a physical blend of controls is licensed only if the model separately shows that an admissible control realizes the corresponding margin vector.
\end{theorem}

The proof (\ref{app:general}) is a separating-hyperplane argument; its content is displayed in Figure~\ref{fig:geometry}. Weight-adaptive acceptance of $s$ means that $s$ is not strictly separable from the up-closed convex hull of the required margins -- for every weighting some plan's margin is cheaper than $s$ in that weighting -- while coordinate-wise acceptance means that $s$ dominates some margin outright. The two coincide exactly when the up-closure of $\mathcal D$ is already convex, in particular whenever one plan dominates all others or the margin menu is convexified. This convex-hull form is shared with scalarization results in multi-objective optimization, but the decision object here is state acceptance under a finite menu. In the witness, each principal plan dips a single coordinate, so the synchronization condition holds trivially and the exact-tube protocol agrees with the margin-first formula. In another tube, scalarizing before taking the path maximum can give a smaller requirement than $w\cdot d_a$, so the convex-hull formula need not be exact. The physical blend in Theorem~\ref{thm:blend} is justified by the witness's linear-flow assumption, not by the abstract theorem alone.

\begin{remark}[fixed-weight countermechanism]\label{rem:fixedweight}
The quantifier gap is sufficient for false acceptance but not necessary. On the witness datum at $s_1=s_2=6/5$ with reserve $x=1/2$, the single fixed weight $w=(1,1)$ already accepts a primitive plan, because its scalar margin meets $s_1+s_2=12/5\ge2$, while no primitive plan meets both typed floors. Enlarging the weight family to the full cone makes the scalarized protocol stricter than any one nonempty subfamily (Proposition~\ref{prop:monotone}), but does not turn this point into a typed-safe one. The assertion concerns the stated action-indexed path/margin protocol of Lemma~\ref{lem:sync}; it is not an identity between a scalarized and a componentwise minimum taken at different times.
\end{remark}

\begin{figure}[!htbp]
\centering
\includegraphics[width=\linewidth]{fig_finite_menu_geometry.png}
\caption{The finite-menu geometric characterisation (Theorem~\ref{thm:menu}) in the plane of two floor margins. Green: coordinate-wise acceptance, the union of the orthants $d_a+\mathbb R^2_+$ (some plan dominates). Black boundary: scalarized acceptance, $\conv(\mathcal D)+\mathbb R^2_+$. Hatched: the acceptance gap, where every weighting licenses some plan but no single plan keeps every floor. (a) A three-plan menu. (b) The margin-class plans of the witness datum, $d_{\textsc{fast}}=(2,0)$ and $d_{\textsc{slow}}=(0,2)$: the gap is the discrepancy triangle $\Q=\{s_1<2,\,s_2<2,\,s_1+s_2\ge2\}$ of Theorem~\ref{thm:witness}, with the witness point $(6/5,6/5)$ in its interior; the reserve-financed plan, which is outside the margin class, serves the triangle only when the reserve covers its cost.}
\label{fig:geometry}
\end{figure}
